% GENERATED by paper/make_tables.py -- do not edit by hand.

\begin{table}[t]
\caption{\textbf{E1: recovery of the analytic label noise.} The target column is
$\mathbb{E}_z[H(p^*)]$, computed in closed form from the generative model, not
estimated. The aleatoric estimate tracks it across the sweep while the epistemic
term stays flat at $O(10^{-4})$ --- both from the same forward passes, so no
rescaling of ``uncertainty'' could make one column follow a moving target while
the other does not. $\Delta$ is positive at \emph{every} point, which the
identifiability argument of Section~\ref{sec:bracket} predicts in advance: a
suboptimal view of the latent can only inflate apparent noise, never deflate it.
The final column is the per-example correlation between the estimated aleatoric
term and the oracle entropy $H(p^*(y|x))$ on the same image, which an estimator
that merely matched the sweep means could not achieve.}
\label{tab:e1}
\centering
\small
\begin{tabular}{cccccccc}
\toprule
$\beta$ & target $\mathbb{E}_z[H(p^*)]$ & est.\ aleatoric & $\Delta$ &
epist.\ ($\times 10^{-3}$) & Bayes err. & model err. & corr.\ (per-ex.) \\
\midrule
0.6 & 0.653 & 0.662 & $0.0087$ & 0.29 & 0.387 & 0.396 & 0.742 \\
1.0 & 0.599 & 0.609 & $0.0102$ & 0.36 & 0.325 & 0.342 & 0.809 \\
1.6 & 0.513 & 0.540 & $0.0269$ & 0.51 & 0.256 & 0.268 & 0.863 \\
2.5 & 0.406 & 0.424 & $0.0181$ & 0.60 & 0.189 & 0.202 & 0.902 \\
4.0 & 0.290 & 0.307 & $0.0167$ & 0.56 & 0.128 & 0.144 & 0.928 \\
\bottomrule
\end{tabular}
\end{table}

\begin{table}[t]
\caption{\textbf{E2: epistemic uncertainty decays with data; aleatoric does not.}
Both columns are computed from the same forward passes, so no monotone rescaling
of ``uncertainty'' could produce a decaying column beside a flat one. Fitted
log--log slope of epistemic against $n$: $-0.85$.}
\label{tab:e2}
\centering
\small
\begin{tabular}{ccccc}
\toprule
$n_{\text{train}}$ & epist.\ ($\times 10^{-3}$) & est.\ aleatoric &
target aleatoric & model err. \\
\midrule
250 & 5.35 & 0.514 & 0.513 & 0.272 \\
500 & 2.73 & 0.496 & 0.513 & 0.267 \\
1000 & 0.52 & 0.529 & 0.513 & 0.272 \\
2000 & 2.32 & 0.504 & 0.513 & 0.275 \\
4000 & 0.31 & 0.528 & 0.513 & 0.265 \\
\bottomrule
\end{tabular}
\end{table}

\begin{table}[t]
\caption{\textbf{E3: calibration before and after temperature scaling}
($T = 1.00$, fitted on in-distribution validation data only).
The correction is a null: every scaled ECE lies inside the raw estimate's $95\%$
bootstrap interval, because a deep ensemble with a heteroscedastic head is
already calibrated. What does vary is calibration under shift, which degrades
monotonically with severity and is worst under the modality shift --- a property
of the model that no in-distribution recalibration addresses. Table~\ref{tab:e3b}
supplies the missing contrast. $^\dagger$ marks semantic shift.}
\label{tab:e3}
\centering
\small
\begin{tabular}{lccccc}
\toprule
& & \multicolumn{2}{c}{ECE} & \multicolumn{2}{c}{NLL} \\
\cmidrule(lr){3-4} \cmidrule(lr){5-6}
split & acc. & raw & scaled & raw & scaled \\
\midrule
val & 0.730 & 0.025 & 0.024 & 0.530 & 0.530 \\
test & 0.733 & 0.017 & 0.019 & 0.530 & 0.531 \\
noise (mild) & 0.738 & 0.022 & 0.023 & 0.509 & 0.509 \\
noise (moderate) & 0.744 & 0.023 & 0.023 & 0.538 & 0.539 \\
noise (severe) & 0.662 & 0.072 & 0.073 & 0.610 & 0.612 \\
blur (mild) & 0.718 & 0.025 & 0.025 & 0.553 & 0.553 \\
blur (moderate) & 0.733 & 0.039 & 0.040 & 0.548 & 0.549 \\
blur (severe) & 0.713 & 0.054 & 0.056 & 0.565 & 0.566 \\
modality shift & 0.669 & 0.077 & 0.078 & 0.606 & 0.608 \\
novel morphology & 0.728 & 0.027 & 0.025 & 0.534 & 0.534 \\
decoupled lesion$^\dagger$ & 0.507 & 0.235 & 0.236 & 0.904 & 0.909 \\
\bottomrule
\end{tabular}
\end{table}

\begin{table}[t]
\caption{\textbf{E3b: the same correction on a model that needs it.} The
baseline is a single network, no heteroscedastic head, dropout off, stopped at
the final epoch rather than at best validation NLL --- each of the main model's
calibration mechanisms removed. Its fitted temperature is
$T = 1.03$ against $T = 1.00$
for the ensemble. Accuracy is bit-identical before and after scaling here, since
a single softmax satisfies the argmax-invariance theorem exactly.}
\label{tab:e3b}
\centering
\small
\begin{tabular}{lcccc}
\toprule
& \multicolumn{2}{c}{single model (ECE)} & \multicolumn{2}{c}{ensemble (ECE)} \\
\cmidrule(lr){2-3} \cmidrule(lr){4-5}
split & raw & scaled & raw & scaled \\
\midrule
in-distribution & 0.016 & 0.012 & 0.017 & 0.019 \\
noise (severe) & 0.153 & 0.149 & 0.072 & 0.073 \\
blur (severe) & 0.013 & 0.015 & 0.054 & 0.056 \\
modality shift & 0.063 & 0.059 & 0.077 & 0.078 \\
novel morphology & 0.036 & 0.032 & 0.027 & 0.025 \\
decoupled lesion$^\dagger$ & 0.241 & 0.237 & 0.235 & 0.236 \\
\bottomrule
\end{tabular}
\end{table}

\begin{table}[t]
\caption{\textbf{E4: OOD detection AUROC.} Best score per column in bold.
\texttt{aleatoric} is a \emph{negative control}: if the decomposition is real it
should not win on semantic shift ($^\dagger$), since a novel input is one the
posterior \emph{disagrees} about rather than one every member agrees is ambiguous.}
\label{tab:e4}
\centering
\small
\begin{tabular}{lccccccccc}
\toprule
score & noise (mild) & noise (moderate) & noise (severe) & blur (mild) & blur (moderate) & blur (severe) & modality shift & novel morphology & decoupled lesion$^\dagger$ \\
\midrule
msp & 0.472 & 0.464 & 0.509 & 0.507 & 0.472 & 0.435 & 0.480 & 0.483 & 0.491 \\
entropy & 0.472 & 0.464 & 0.509 & 0.507 & 0.472 & 0.435 & 0.480 & 0.483 & 0.491 \\
energy & 0.492 & 0.507 & 0.520 & \textbf{0.510} & 0.483 & 0.460 & 0.494 & 0.483 & 0.491 \\
epistemic & 0.558 & 0.870 & 0.937 & 0.493 & 0.439 & 0.681 & 0.584 & 0.533 & 0.522 \\
aleatoric & 0.471 & 0.455 & 0.435 & 0.507 & 0.473 & 0.435 & 0.479 & 0.483 & 0.490 \\
mahalanobis & \textbf{0.784} & \textbf{0.986} & \textbf{1.000} & 0.476 & \textbf{0.607} & \textbf{0.917} & \textbf{0.859} & \textbf{0.562} & \textbf{0.548} \\
\bottomrule
\end{tabular}
\end{table}

\begin{table}[t]
\caption{\textbf{E5: selective prediction.} E-AURC isolates ranking quality from
classifier quality. On the mixed stream every novel case counts as an error,
modelling a screening queue that contains scans the model should decline.}
\label{tab:e5}
\centering
\small
\begin{tabular}{lcccc}
\toprule
confidence function & AURC & E-AURC & risk@50\% cov. & cov.@5\% risk \\
\midrule
\multicolumn{5}{l}{\emph{In-distribution}} \\
\quad msp & 0.1489 & 0.1093 & 0.147 & 0.138 \\
\quad neg\_total\_entropy & 0.1489 & 0.1093 & 0.147 & 0.138 \\
\quad neg\_aleatoric & 0.1489 & 0.1093 & 0.147 & 0.138 \\
\quad neg\_epistemic & 0.2461 & 0.2065 & 0.240 & 0.000 \\
\multicolumn{5}{l}{\emph{Mixed ID + novel}} \\
\quad msp & 0.4244 & 0.2674 & 0.425 & 0.001 \\
\quad neg\_total\_entropy & 0.4244 & 0.2674 & 0.425 & 0.001 \\
\quad neg\_epistemic & 0.4955 & 0.3384 & 0.484 & 0.000 \\
\bottomrule
\end{tabular}
\end{table}

\begin{table}[t]
\caption{\textbf{E6: Monte-Carlo budget.} The total term plugs a $T$-sample mean
into the concave functional $H$, so by Jensen's inequality it is under-estimated
at small $T$, and the epistemic term with it. The estimate rises and flattens;
quoting an epistemic number without this curve quotes an unknown fraction of it.}
\label{tab:e6}
\centering
\small
\begin{tabular}{cccc}
\toprule
$T$ & epistemic & aleatoric & total \\
\midrule
2 & 0.00355 & 0.53275 & 0.53631 \\
4 & 0.00544 & 0.53434 & 0.53978 \\
8 & 0.00644 & 0.53421 & 0.54065 \\
16 & 0.00684 & 0.53372 & 0.54056 \\
32 & 0.00709 & 0.53379 & 0.54089 \\
64 & 0.00722 & 0.53371 & 0.54092 \\
\bottomrule
\end{tabular}
\end{table}
